% Options for packages loaded elsewhere
\PassOptionsToPackage{unicode}{hyperref}
\PassOptionsToPackage{hyphens}{url}
\documentclass[
  12pt,
]{ctexart}
\usepackage{xcolor}
\usepackage{amsmath,amssymb}
\setcounter{secnumdepth}{-\maxdimen} % remove section numbering
\usepackage{iftex}
\ifPDFTeX
  \usepackage[T1]{fontenc}
  \usepackage[utf8]{inputenc}
  \usepackage{textcomp} % provide euro and other symbols
\else % if luatex or xetex
  \usepackage{unicode-math} % this also loads fontspec
  \defaultfontfeatures{Scale=MatchLowercase}
  \defaultfontfeatures[\rmfamily]{Ligatures=TeX,Scale=1}
\fi
\usepackage{lmodern}
\ifPDFTeX\else
  % xetex/luatex font selection
\fi
% Use upquote if available, for straight quotes in verbatim environments
\IfFileExists{upquote.sty}{\usepackage{upquote}}{}
\IfFileExists{microtype.sty}{% use microtype if available
  \usepackage[]{microtype}
  \UseMicrotypeSet[protrusion]{basicmath} % disable protrusion for tt fonts
}{}
\makeatletter
\@ifundefined{KOMAClassName}{% if non-KOMA class
  \IfFileExists{parskip.sty}{%
    \usepackage{parskip}
  }{% else
    \setlength{\parindent}{0pt}
    \setlength{\parskip}{6pt plus 2pt minus 1pt}}
}{% if KOMA class
  \KOMAoptions{parskip=half}}
\makeatother
\usepackage{longtable,booktabs,array}
\usepackage{calc} % for calculating minipage widths
% Correct order of tables after \paragraph or \subparagraph
\usepackage{etoolbox}
\makeatletter
\patchcmd\longtable{\par}{\if@noskipsec\mbox{}\fi\par}{}{}
\makeatother
% Allow footnotes in longtable head/foot
\IfFileExists{footnotehyper.sty}{\usepackage{footnotehyper}}{\usepackage{footnote}}
\makesavenoteenv{longtable}
\setlength{\emergencystretch}{3em} % prevent overfull lines
\providecommand{\tightlist}{%
  \setlength{\itemsep}{0pt}\setlength{\parskip}{0pt}}
\usepackage{bookmark}
\IfFileExists{xurl.sty}{\usepackage{xurl}}{} % add URL line breaks if available
\urlstyle{same}
\hypersetup{
  hidelinks,
  pdfcreator={LaTeX via pandoc}}

\author{SCX}
\date{2026}

\begin{document}

\section{SCX
知识产权保护、开源策略与商业模式}\label{scx-ux77e5ux8bc6ux4ea7ux6743ux4fddux62a4ux5f00ux6e90ux7b56ux7565ux4e0eux5546ux4e1aux6a21ux5f0f}

\begin{quote}
整理自 GPT 讨论 (2026-06-26)，综合 IP
保护、开源发布、商业化和组织架构设计。
\end{quote}

\begin{center}\rule{0.5\linewidth}{0.5pt}\end{center}

\subsection{一、五层资产模型}\label{ux4e00ux4e94ux5c42ux8d44ux4ea7ux6a21ux578b}

将 SCX 资产分为五层，每层采用不同的发表/开源/保留策略：

\begin{longtable}[]{@{}
  >{\raggedright\arraybackslash}p{(\linewidth - 4\tabcolsep) * \real{0.3333}}
  >{\raggedright\arraybackslash}p{(\linewidth - 4\tabcolsep) * \real{0.3333}}
  >{\raggedright\arraybackslash}p{(\linewidth - 4\tabcolsep) * \real{0.3333}}@{}}
\toprule\noalign{}
\begin{minipage}[b]{\linewidth}\raggedright
层级
\end{minipage} & \begin{minipage}[b]{\linewidth}\raggedright
内容
\end{minipage} & \begin{minipage}[b]{\linewidth}\raggedright
策略
\end{minipage} \\
\midrule\noalign{}
\endhead
\bottomrule\noalign{}
\endlastfoot
\textbf{数学理论} & SCX 定义、状态条件风险、数据价值函数、命题证明 &
\textbf{发表，完全公开} \\
\textbf{算法框架} & pseudocode、核心流程、基本公式 &
\textbf{发表，完全公开} \\
\textbf{最小代码} & toy demo、基础 SCX score、简单 coreset &
\textbf{开源} \\
\textbf{工程实现} & 高性能
pipeline、自动诊断、可视化、工业接口、模型适配器 &
\textbf{暂不全开源} \\
\textbf{数据/模型/专家库} & 高质量
benchmarks、专家势函数库、私有筛选规则、校准参数 &
\textbf{重点商业资产，严格保留} \\
\end{longtable}

\subsubsection{做什么}\label{ux505aux4ec0ux4e48}

\begin{itemize}
\tightlist
\item
  论文里把数学定义、命题、伪代码、synthetic 实验全部公开。
\item
  开源一个轻量 \texttt{scx-core}，让别人能复现论文结论。
\item
  工程实现可以开源部分通用模块，但核心策略、阈值、校准参数保留。
\end{itemize}

\subsubsection{不做什么}\label{ux4e0dux505aux4ec0ux4e48}

\begin{itemize}
\tightlist
\item
  不把所有生产级代码、数据、模型和调参细节一次性放出去。
\item
  不把工业级 pipeline、商业 GUI、专家库、训练好的高价值模型开源。
\item
  不把私有校准参数和阈值策略公开。
\end{itemize}

\begin{center}\rule{0.5\linewidth}{0.5pt}\end{center}

\subsection{二、三层开源策略}\label{ux4e8cux4e09ux5c42ux5f00ux6e90ux7b56ux7565}

\subsubsection{第一层：论文完全公开}\label{ux7b2cux4e00ux5c42ux8bbaux6587ux5b8cux5168ux516cux5f00}

\textbf{做什么：} - 公开 SCX
核心数学定义、状态条件专家风险、数据价值函数、冗余压缩目标、专家路由公式。
- 公开命题与证明、合成实验、小规模真实案例、算法伪代码。 -
目标期刊/会议：Nature Computational Science / npj Digital Medicine /
TMLR / JMLR 等。

\textbf{目的：} 抢概念、拿引用、建立信誉、获得学术优先权。

\textbf{不做什么：} -
论文里不要公开全部工程细节、自动阈值选择、高性能状态划分策略。 -
不要把大规模工业 pipeline 写进论文补充材料。

\begin{center}\rule{0.5\linewidth}{0.5pt}\end{center}

\subsubsection{第二层：reference code
开源}\label{ux7b2cux4e8cux5c42reference-code-ux5f00ux6e90}

\textbf{做什么：} - 开源轻量仓库 \texttt{scx-core}，包含：state
partition、expert reliability estimation、redundancy compression、noise
score、data value score、toy examples。 - 许可证偏宽松（如 MIT / Apache
2.0），方便学术传播。 - 开源 \texttt{SCX-Health} 仓库，用 MedMNIST /
HAM10000 / CheXpert 做公开 demo。

\textbf{目的：} 让别人能复现论文、建立开源社区、吸引合作者、获得公信力。

\textbf{不做什么：} - 不开源工业级数据处理 pipeline。 - 不开源企业级
API、商业 dashboard、自动报告系统。 - 不开源完整 AlGaN/NEP/MACE
专家库和大规模 benchmark 原始数据。

\subsubsection{第三层：商业版闭源}\label{ux7b2cux4e09ux5c42ux5546ux4e1aux7248ux95edux6e90}

\textbf{SCX-Pro / SCX-Compiler 保留内容：} - 大规模数据管线 - MLIP
expert compiler（能量零点对齐、gauge normalization、domain
certificate、conflict arbitration、distilled student） - LLM data
curation 模块 - 企业可视化 dashboard - 自动报告生成 - 大模型数据筛选模块
- 云 API 和企业私有部署 - 私有 benchmark

\textbf{目的：} 服务商业用户，形成付费壁垒。

\textbf{不做什么：} - 不把商业版放在公开仓库。 -
不把企业部署代码在开源协议下发布。

\subsubsection{核心原则}\label{ux6838ux5fc3ux539fux5219}

\begin{quote}
\textbf{开源到''足以复现论文''，不要开源到''足以替代你的商业产品''。}
\end{quote}

\begin{center}\rule{0.5\linewidth}{0.5pt}\end{center}

\subsection{三、IP 保护流程}\label{ux4e09ip-ux4fddux62a4ux6d41ux7a0b}

\subsubsection{论文发表前必须做的事}\label{ux8bbaux6587ux53d1ux8868ux524dux5fc5ux987bux505aux7684ux4e8b}

正确顺序：

\begin{verbatim}
内部 invention disclosure → 专利评估 → 问学校/专利代理人
→ 值得就提交申请 → arXiv/投稿 → 选择性开源 → 商业版
\end{verbatim}

\textbf{做什么：} 1. \textbf{写 invention disclosure（内部 2-5
页材料）}：记录 SCX
核心流程、数据压缩、专家路由、状态价值函数、应用场景。这是确权第一步。
2.
\textbf{找学校知识产权办公室或校外专利代理人评估}：确认哪些有专利价值。
3. \textbf{有专利价值就先申请}：至少把关键权利要求占住，再 arXiv/投稿。
4. \textbf{论文发表后再开源最小代码}：确保公开披露不破坏新颖性。 5.
\textbf{最后做商业版}。

\textbf{不做什么：} -
不要先公开完整论文和代码后再想专利------许多司法辖区的宽限期有限，不值得赌。
- 不要在论文/arXiv/代码公开前在朋友圈、报告、PPT 中披露关键技术方案。

\subsubsection{哪些可以专利？}\label{ux54eaux4e9bux53efux4ee5ux4e13ux5229}

这些具体技术方案可以申请专利或软著： - 状态条件数据价值评估系统 -
SCX-Compress 冗余压缩流程（具体技术流程） - 专家可靠性矩阵与数据路由策略
- 高误差状态 acquire/relabel/downweight 策略 - 用于 MLIP/LLM
数据审计的自动报告系统

\subsubsection{哪些靠论文确权？}\label{ux54eaux4e9bux9760ux8bbaux6587ux786eux6743}

这些靠论文建立学术优先权（专利不适合保护纯数学）： - 数学公式本身 -
抽象思想、框架概念 - 定义、命题、证明

\subsubsection{``防御性公开''策略}\label{ux9632ux5fa1ux6027ux516cux5f00ux7b56ux7565}

不是把全部资产藏起来，而是：

\begin{verbatim}
论文（确立学术优先权）
+ 选择性开源（建立生态）
+ 闭源增强版（保护商业核心）
\end{verbatim}

三者并行，形成 ``公开理论，保留产品化能力'' 的防御体系。

\begin{center}\rule{0.5\linewidth}{0.5pt}\end{center}

\subsection{四、学校 IP
风险切割}\label{ux56dbux5b66ux6821-ip-ux98ceux9669ux5207ux5272}

\subsubsection{主要风险}\label{ux4e3bux8981ux98ceux9669}

如果使用学校资源（超算、导师项目、实验室数据），成果可能被认定为职务发明/职务作品。

\textbf{高风险：} DFT 数据、AlGaN/MLIP 应用软件、用学校超算产出的结果。
\textbf{可切割：} SCX
通用数学框架与一般代码（如果独立开发、使用非学校资源）。

\subsubsection{做什么}\label{ux505aux4ec0ux4e48-1}

\begin{enumerate}
\def\labelenumi{\arabic{enumi}.}
\tightlist
\item
  \textbf{新建私人仓库}：用个人邮箱、个人 GitHub、个人电脑，建立
  \texttt{scx-core-private}。
\item
  \textbf{开启可验证时间线}：Git commit 全部保留，关键文档生成 SHA-256
  哈希，定期做时间戳存证，保留电脑购买记录。
\item
  \textbf{Clean-room 重写}：SCX 通用版重新写一遍，不从旧项目复制粘贴。
\item
  \textbf{写
  DEVELOPMENT\_LOG.md}：记录开发日期、地点、设备、是否使用学校资源。
\item
  \textbf{把 DFT/AlGaN 和 SCX 分离}：SCX 论文不使用学校 DFT
  数据作为核心证据。
\item
  \textbf{查清签过的协议}：入学协议、奖学金协议、课题组管理规定、学校 IP
  政策、超算使用协议。
\end{enumerate}

\subsubsection{不做什么}\label{ux4e0dux505aux4ec0ux4e48-1}

\begin{itemize}
\tightlist
\item
  不把 SCX 核心代码上传学校服务器。
\item
  不在学校超算跑 SCX。
\item
  不用导师项目数据做 SCX 核心实验。
\item
  不用课题组网盘备份核心代码。
\item
  不在课题组群里发完整设计。
\item
  不把 repo 权限给课题组任何人。
\item
  不要轻易把导师列为 SCX 数学论文的共同作者（如果导师没有实质参与）。
\end{itemize}

\begin{center}\rule{0.5\linewidth}{0.5pt}\end{center}

\subsection{五、商业模式}\label{ux4e94ux5546ux4e1aux6a21ux5f0f}

\subsubsection{四种产品化方向}\label{ux56dbux79cdux4ea7ux54c1ux5316ux65b9ux5411}

\paragraph{1. 数据压缩 /
数据去冗余工具（SCX-Compress）}\label{ux6570ux636eux538bux7f29-ux6570ux636eux53bbux5197ux4f59ux5de5ux5177scx-compress}

\begin{itemize}
\tightlist
\item
  用户输入训练数据，输出：保留多少数据、训练误差保持多少、高误差区域不删、边界
  anchor 保留、噪声点降权。
\item
  最直接的商业价值。
\end{itemize}

\paragraph{2.
专家标注路由工具}\label{ux4e13ux5bb6ux6807ux6ce8ux8defux7531ux5de5ux5177}

\begin{itemize}
\tightlist
\item
  适用于大模型、科学数据、人类专家、DFT 标注等场景。
\item
  判断：哪类数据交给哪个 expert、哪类数据必须高精度
  oracle、哪类数据不用标（冗余）。
\item
  服务于 AI 数据公司、材料计算团队、自动驾驶/医疗影像/遥感等领域。
\end{itemize}

\paragraph{3. MLIP expert compiler}\label{mlip-expert-compiler}

\begin{itemize}
\tightlist
\item
  最自然的垂直商业入口。
\item
  输入：多个 ACE/NEP/MACE expert → 能量零点对齐 → gauge normalization →
  domain certificate → conflict arbitration → distilled student / merged
  potential。
\end{itemize}

\paragraph{4. 数据价值审计报告（SCX Data Value Audit
Report）}\label{ux6570ux636eux4ef7ux503cux5ba1ux8ba1ux62a5ux544ascx-data-value-audit-report}

\begin{itemize}
\tightlist
\item
  交付物包括：数据冗余地图、噪声风险地图、高价值状态地图、专家可靠性矩阵、建议保留/删除/重标注/补数据清单、预计压缩比例、压缩前后模型性能对比、可复现实验脚本。
\item
  是最容易商业化的形式------企业愿意为''审计能力''买单。
\end{itemize}

\subsubsection{五种收费模式}\label{ux4e94ux79cdux6536ux8d39ux6a21ux5f0f}

\paragraph{模式 A：赞助研究（Sponsored
Research）}\label{ux6a21ux5f0f-aux8d5eux52a9ux7814ux7a76sponsored-research}

\begin{itemize}
\tightlist
\item
  公司给课题组一笔钱，支持 SCX 在其场景上验证。
\item
  交付：技术报告、论文合作、内部 demo、数据审计结果。
\item
  适合博士身份，缺点是学校可能参与管理经费和 IP。
\end{itemize}

\paragraph{模式 B：产业联盟会员（SCX
Consortium）}\label{ux6a21ux5f0f-bux4ea7ux4e1aux8054ux76dfux4f1aux5458scx-consortium}

\begin{itemize}
\tightlist
\item
  会员公司每年交会费，共同支持 SCX 标准。
\item
  每家公司私有数据不共享，共同获得工具更新、benchmark、报告模板。
\item
  最符合''他们任何一个人都不希望我被对方收购''的中立性要求。
\end{itemize}

\paragraph{模式
C：第三方审计服务}\label{ux6a21ux5f0f-cux7b2cux4e09ux65b9ux5ba1ux8ba1ux670dux52a1}

\begin{itemize}
\tightlist
\item
  公司给数据，你出审计报告。
\item
  报告包括：Dataset redundancy score、Noise risk score、Valuable state
  map、Expert reliability matrix、Recommended compression
  ratio、Recommended relabeling set。
\item
  最容易商业化的形式。
\end{itemize}

\paragraph{模式
D：非独占授权}\label{ux6a21ux5f0f-dux975eux72ecux5360ux6388ux6743}

\begin{itemize}
\tightlist
\item
  把 SCX-Pro 授权给多家公司使用。
\item
  \textbf{关键：非独占授权，不卖断，不给某一家独占。}
\end{itemize}

\paragraph{模式 E：开源核心 + 企业版（Open
Core）}\label{ux6a21ux5f0f-eux5f00ux6e90ux6838ux5fc3-ux4f01ux4e1aux7248open-core}

\begin{itemize}
\tightlist
\item
  开源 SCX-Core。
\item
  收费
  SCX-Pro：大规模数据处理、自动报告、多专家模型接口、企业隐私部署、高性能压缩、可视化
  dashboard。
\end{itemize}

\subsubsection{``中立第三方数据价值审计方''的定位}\label{ux4e2dux7acbux7b2cux4e09ux65b9ux6570ux636eux4ef7ux503cux5ba1ux8ba1ux65b9ux7684ux5b9aux4f4d}

核心定位： \textgreater{} \textbf{SCX 是面向 AI
训练流程的第三方数据价值审计框架。}

不卖''一个公式''，而是卖： - Data Value Audit / Expert Reliability Audit
/ Dataset Compression Report -
``帮我减少训练数据量或标注成本，精度基本不掉''

\subsubsection{与公司谈话术}\label{ux4e0eux516cux53f8ux8c08ux8bddux672f}

正确说法： \textgreater{}
我不做独占外包，也不做单家公司内部工具。我做的是中立的 SCX
数据价值评估标准。贵公司可以作为早期产业合作方，提供资金、场景、工程师参与共建；你们获得优先试用、内部部署和非独占授权，但核心标准和
IP 保持中立，以确保其他公司也能信任这个体系。

更商业的版本： \textgreater{} SCX
不出售给单一公司。我们提供非独占授权、数据审计、联合 benchmark
和企业部署。所有成员公司都可以使用，但没有一家可以控制标准。

\begin{center}\rule{0.5\linewidth}{0.5pt}\end{center}

\subsection{六、组织架构}\label{ux516dux7ec4ux7ec7ux67b6ux6784}

\subsubsection{三层结构}\label{ux4e09ux5c42ux7ed3ux6784}

\paragraph{第一层：SCX
Research（学术层）}\label{ux7b2cux4e00ux5c42scx-researchux5b66ux672fux5c42}

\begin{itemize}
\tightlist
\item
  负责：论文、数学定义、benchmark、reference
  implementation、开源核心、学术影响力。
\item
  挂靠形式：学校、实验室、个人开源项目、非营利式研究组织。
\item
  控制人：你（scientific founder / chief scientist）。
\end{itemize}

\paragraph{第二层：SCX
Consortium（产业联盟层）}\label{ux7b2cux4e8cux5c42scx-consortiumux4ea7ux4e1aux8054ux76dfux5c42}

\begin{itemize}
\tightlist
\item
  成员公司交钱，获得：标准制定参与权、内部试用权、benchmark
  访问权、定期技术报告、优先试点机会、非独占企业授权折扣。
\item
  成员不能获得的：SCX 核心 IP
  所有权、独占授权、单方面修改标准的权力、收购后封锁其他成员的权力。
\end{itemize}

\paragraph{第三层：SCX Operator /
SCX-Pro（商业执行层）}\label{ux7b2cux4e09ux5c42scx-operator-scx-proux5546ux4e1aux6267ux884cux5c42}

\begin{itemize}
\tightlist
\item
  负责：企业项目交付、数据审计、工程部署、可视化平台、API、客户支持、合同和收费。
\item
  可由 CEO/COO/项目经理/工程负责人管理，你不需要天天打商业仗。
\item
  你保留：scientific founder、chief scientist、standards
  chair、architecture veto、IP owner/licensor。
\end{itemize}

\subsubsection{你应控制的四项核心权力}\label{ux4f60ux5e94ux63a7ux5236ux7684ux56dbux9879ux6838ux5fc3ux6743ux529b}

\begin{enumerate}
\def\labelenumi{\arabic{enumi}.}
\tightlist
\item
  \textbf{IP 控制权}：SCX 名字、核心框架、代码主仓库、商标、专利/软著。
\item
  \textbf{标准控制权}：什么叫 SCX-compatible、什么叫
  SCX-certified，由你定义。
\item
  \textbf{认证控制权}：企业想说''我通过 SCX
  数据价值审计''，需要你的体系认证。
\item
  \textbf{方向控制权}：哪些模块进入
  core、哪些只是企业插件，你有最终路线权。
\end{enumerate}

\subsubsection{你不需要控制的}\label{ux4f60ux4e0dux9700ux8981ux63a7ux5236ux7684}

\begin{itemize}
\tightlist
\item
  每个客户怎么谈
\item
  每个工程需求怎么排
\item
  每个 dashboard 怎么做
\item
  每个部署细节
\end{itemize}

这些交给运营团队。

\begin{center}\rule{0.5\linewidth}{0.5pt}\end{center}

\subsection{七、与公司合作的注意事项}\label{ux4e03ux4e0eux516cux53f8ux5408ux4f5cux7684ux6ce8ux610fux4e8bux9879}

\subsubsection{做什么}\label{ux505aux4ec0ux4e48-2}

\begin{enumerate}
\def\labelenumi{\arabic{enumi}.}
\tightlist
\item
  \textbf{合同写成 sponsored research / membership / paid pilot}，而不是
  equity acquisition / exclusive buyout。
\item
  \textbf{坚持非独占授权，不卖断。}
\item
  \textbf{甲方出人可以，但代码贡献必须受控：}

  \begin{itemize}
  \tightlist
  \item
    签 Contributor License Agreement (CLA)。
  \item
    贡献进入 SCX core 前必须签 CLA。
  \item
    SCX 维护方有最终 merge 权。
  \item
    贡献不得附带独占限制。
  \item
    贡献后的 core 仍由 SCX 标准方统一维护。
  \end{itemize}
\item
  \textbf{技术路线你有最终权}：

  \begin{itemize}
  \tightlist
  \item
    标准命名权、版本发布权、reference implementation 合并权、benchmark
    认证口径、SCX 商标使用权、论文署名与发布权。
  \item
    公司可以投票建议，但不能控制核心标准。
  \end{itemize}
\item
  \textbf{找可靠的小班子}：运营负责人、工程负责人、学术/算法核心（你+1-2学生）、法务/IP
  顾问、商务顾问。
\end{enumerate}

\subsubsection{不做什么}\label{ux4e0dux505aux4ec0ux4e48-2}

\begin{enumerate}
\def\labelenumi{\arabic{enumi}.}
\tightlist
\item
  \textbf{不签独占}：比如''公司 A 出钱，你把 SCX 在 AI
  数据领域独占授权给 A''------这直接毁掉中立性。
\item
  \textbf{不接受''买断式合作开发''}：比如''你帮我们做，成果归我们''。
\item
  \textbf{不让甲方控制论文发表}：合作协议里可以有保密审查期，但不要给甲方无限否决发表的权力。
\item
  \textbf{不让甲方工程师绕过主仓库}：所有核心贡献必须进统一代码库、统一协议、统一
  review。
\item
  \textbf{不要私下收大额商业款}：你是博士，学校 IP 和经费合规要处理好。
\item
  \textbf{不要过早卖给一家大公司}：如果早早被收购，其他公司不会再用你的标准。
\item
  \textbf{不要全开源生产级代码}：研究版开源，工业版闭源。
\end{enumerate}

\subsubsection{如何让甲方出钱出人但不失控}\label{ux5982ux4f55ux8ba9ux7532ux65b9ux51faux94b1ux51faux4ebaux4f46ux4e0dux5931ux63a7}

关键设计：

\begin{verbatim}
公司 A 出钱 + 工程师
公司 B 出钱 + 数据案例
公司 C 出钱 + 应用场景
         ↓
    SCX 中立实验室/标准方
         ↓
统一标准、开源核心、认证报告、商业审计、企业部署
\end{verbatim}

你不是给某一家打工，而是让多方围绕你的标准参与。

\begin{center}\rule{0.5\linewidth}{0.5pt}\end{center}

\subsection{八、商业路线图}\label{ux516bux5546ux4e1aux8defux7ebfux56fe}

\subsubsection{阶段 1：论文确权 +
最小工具（现在）}\label{ux9636ux6bb5-1ux8bbaux6587ux786eux6743-ux6700ux5c0fux5de5ux5177ux73b0ux5728}

\begin{itemize}
\tightlist
\item
  SCX 数学论文（arXiv 或投稿）。
\item
  SCX-Core 最小代码开源。
\item
  一个漂亮 demo（MedMNIST 数据压缩）。
\item
  一个数据压缩/冗余案例。
\end{itemize}

\subsubsection{阶段 2：找 1-2 个甲方做 paid
pilot}\label{ux9636ux6bb5-2ux627e-1-2-ux4e2aux7532ux65b9ux505a-paid-pilot}

\begin{itemize}
\tightlist
\item
  每个 pilot 只做一件事：帮它减少训练数据量或标注成本。
\item
  交付报告，不承诺全平台。
\end{itemize}

\subsubsection{阶段 3：成立 SCX Working Group /
Consortium}\label{ux9636ux6bb5-3ux6210ux7acb-scx-working-group-consortium}

\begin{itemize}
\tightlist
\item
  让多个单位加入：公司、高校、材料计算团队、AI 数据团队。
\item
  收 membership fee 或项目费。
\end{itemize}

\subsubsection{阶段
4：找职业经理人}\label{ux9636ux6bb5-4ux627eux804cux4e1aux7ecfux7406ux4eba}

\begin{itemize}
\tightlist
\item
  你做 scientific founder，别人做 CEO/COO。
\item
  你保留标准控制权和技术 veto。
\end{itemize}

\subsubsection{阶段
5：建立认证体系}\label{ux9636ux6bb5-5ux5efaux7acbux8ba4ux8bc1ux4f53ux7cfb}

\begin{itemize}
\tightlist
\item
  SCX-certified dataset
\item
  SCX-compressed dataset
\item
  SCX expert reliability report
\item
  SCX-compatible data pipeline
\end{itemize}

\subsubsection{明确不能做的事}\label{ux660eux786eux4e0dux80fdux505aux7684ux4e8b}

\begin{quote}
\begin{itemize}
\tightlist
\item
  不要一上来找大厂签''所有成果归甲方''的合作。
\item
  不要把 SCX 讲成万能 AI，要讲成''状态条件质量判断层''。
\item
  不要喊''取代 TCAD/OPC/Calibre''，要说''SCX complements physical
  solvers by deciding when their outputs are trustworthy''。
\item
  LLM 可以放 discussion，主战场是高精度标签/仿真数据。
\item
  不要一开始写成商业帝国 PPT------先做出第一个''能省 50\%
  仿真/标注成本且性能不掉''的样板。
\end{itemize}
\end{quote}

\begin{center}\rule{0.5\linewidth}{0.5pt}\end{center}

\subsection{九、论文谱系与发表策略}\label{ux4e5dux8bbaux6587ux8c31ux7cfbux4e0eux53d1ux8868ux7b56ux7565}

\subsubsection{推荐发表顺序}\label{ux63a8ux8350ux53d1ux8868ux987aux5e8f}

\begin{longtable}[]{@{}
  >{\raggedright\arraybackslash}p{(\linewidth - 6\tabcolsep) * \real{0.2222}}
  >{\raggedright\arraybackslash}p{(\linewidth - 6\tabcolsep) * \real{0.2222}}
  >{\raggedright\arraybackslash}p{(\linewidth - 6\tabcolsep) * \real{0.2222}}
  >{\raggedright\arraybackslash}p{(\linewidth - 6\tabcolsep) * \real{0.3333}}@{}}
\toprule\noalign{}
\begin{minipage}[b]{\linewidth}\raggedright
顺序
\end{minipage} & \begin{minipage}[b]{\linewidth}\raggedright
论文
\end{minipage} & \begin{minipage}[b]{\linewidth}\raggedright
定位
\end{minipage} & \begin{minipage}[b]{\linewidth}\raggedright
目标期刊
\end{minipage} \\
\midrule\noalign{}
\endhead
\bottomrule\noalign{}
\endlastfoot
1 & \textbf{EGP Paper 1}: ACE gauge-normalized expert merging &
材料势函数方法，根据地（有完整 DFT 数据） & npj Computational Materials
/ PRM / JCTC \\
2 & \textbf{SCX-Theory}: 数学定义、命题、可识别性边界 & 理论地基，arXiv
占坑 & TMLR / SIMODS / JMLR \\
3 & \textbf{SCX-MLIP}: SCX 理论应用于 MLIP & 理论→应用闭环 & Nature
Communications / npj Computational Materials \\
4 & \textbf{SCX-Sim}: 科学与工程仿真多保真调度 & 跨领域大文章，主冲
Nature 系列 & Nature Computational Science \\
5 & \textbf{SCX-Health}: 医学/视觉数据审计 & 开源影响力 + 商业入口 & npj
Digital Medicine \\
\end{longtable}

\subsubsection{每篇的独立性要求}\label{ux6bcfux7bc7ux7684ux72ecux7acbux6027ux8981ux6c42}

\begin{itemize}
\tightlist
\item
  \textbf{EGP Paper 1} 解决 ACE expert 的 gauge/energy
  规范化与合并方法。
\item
  \textbf{SCX-Theory}
  提出状态条件专家性与数据价值的数学定义和可识别性边界。
\item
  \textbf{SCX-MLIP} 将 SCX 理论应用回 MLIP 领域，完成理论→应用闭环。
\item
  \textbf{SCX-Sim} 把 SCX 发展成高保真科学/工程仿真的多保真调度系统。
\item
  \textbf{SCX-Health}
  解决医学数据中的标签噪声、长尾、专家复核和冗余压缩。
\end{itemize}

不要把 SCX 当成一个方法拆成很多应用论文；要把 SCX
当成一个理论母体，每篇论文解决一个不同层级的问题。

\end{document}
